% !TEX program = xelatex
% 使用 XeLaTeX 编译

\documentclass[12pt, a4paper]{article}

% ---------- 字体与编码 ----------
\usepackage{fontspec}
\usepackage{xeCJK}
\setmainfont{Times New Roman}
\setCJKmainfont{SimSun} % 确保编译环境有宋体

% ---------- 数学宏包 ----------
\usepackage{amsmath, amssymb, amsthm}
\usepackage{mathrsfs} % 花体字母

% ---------- 页面边距 ----------
\usepackage{geometry}
\geometry{left=3cm, right=3cm, top=2.5cm, bottom=2.5cm}

% ---------- 定理环境定义 ----------
\newtheorem{definition}{定义}[section]
\newtheorem{theorem}{定理}[section]
\newtheorem{lemma}{引理}[section]
\newtheorem{corollary}{推论}[section]
\newtheorem{example}{例}[section]

% ---------- 自定义命令 ----------
\newcommand{\R}{\mathbb{R}} % 实数集
\newcommand{\N}{\mathbb{N}} % 自然数集
\newcommand{\limn}{\lim_{n \to \infty}} % 极限简写

% ---------- 文档开始 ----------
\begin{document}

\title{数学分析 II：极限与连续}
\author{inshergorohl}
\date{Flux 1.3 Reference}
\maketitle

\section{数列极限}
\subsection{定义与性质}

\begin{definition}[$\varepsilon$-N 语言]
设 $\{x_n\}$ 为数列，$a$ 为定数。若 $\forall \varepsilon > 0$，$\exists N \in \N$，使得当 $n > N$ 时，有
\[
|x_n - a| < \varepsilon,
\]
则称数列 $\{x_n\}$ \textbf{收敛}于 $a$，记作
\[
\limn x_n = a.
\]
\end{definition}

\begin{theorem}[极限的唯一性]
若数列 $\{x_n\}$ 收敛，则其极限唯一。
\end{theorem}

\begin{proof}
反证法。假设 $\limn x_n = a$ 且 $\limn x_n = b$，且 $a \neq b$。
取 $\varepsilon = \frac{|a-b|}{2} > 0$。
由极限定义，$\exists N_1, N_2$，使得...
（此处省略证明细节，保持简洁）
\end{proof}

\section{函数极限}

\begin{definition}[单侧极限]
设函数 $f(x)$ 在 $x_0$ 的某个去心邻域内有定义。
\[
\lim_{x \to x_0^+} f(x) = A \iff \forall \varepsilon > 0, \exists \delta > 0, \text{s.t. } 0 < x - x_0 < \delta \implies |f(x)-A|<\varepsilon.
\]
\end{definition}

\section{连续性}

\begin{definition}[连续]
函数 $f(x)$ 在点 $x_0$ 处连续，当且仅当：
\[
\lim_{x \to x_0} f(x) = f(x_0).
\]
\end{definition}

\begin{theorem}[介值定理]
设 $f \in C[a,b]$，且 $f(a) \neq f(b)$。则 $\forall \mu$ 介于 $f(a)$ 与 $f(b)$ 之间，$\exists \xi \in (a,b)$，使得 $f(\xi) = \mu$。
\end{theorem}

\section*{附录：常用等价无穷小}
\[
\begin{aligned}
&\sin x \sim x, \quad \tan x \sim x, \quad 1 - \cos x \sim \frac{x^2}{2}, \\
&\ln(1+x) \sim x, \quad e^x - 1 \sim x, \quad \sqrt[n]{1+x} - 1 \sim \frac{x}{n}.
\end{aligned}
\]

\end{document}
